% Vietnamese translation for OI Public Library
% Source-PDF: 动态规划 DYNAMIC PROGRAMMING/动态规划的优化 - 唐文斌.pdf
\documentclass[11pt]{article}
\usepackage{oipl-vn}

\title{Tối ưu hóa quy hoạch động}
\author{Đường Văn Bân}
\date{}

\begin{document}
\maketitle

\origtitle{Dynamic Programming Optimization}
\sourcepath{Dynamic Programming / DP Optimization - Tang Wenbin.pdf}

\section*{Tư tưởng chung}

Độ phức tạp thời gian của một lời giải quy hoạch động thường có thể tách
thành hai phần:
\[
  \text{số trạng thái}\times\text{thời gian chuyển cho mỗi trạng thái}.
\]
Vì vậy, khi một quy hoạch động trực tiếp chưa đủ nhanh, ta thường đi theo một
trong hai hướng:
\begin{itemize}
  \item giảm số trạng thái cần xét;
  \item tăng hiệu quả của phép chuyển trạng thái.
\end{itemize}
Tài liệu gốc chủ yếu minh họa hướng thứ hai: giữ mô hình trạng thái tự nhiên,
nhưng nhận ra phép cực trị trong công thức chuyển có thể được thay bằng truy
vấn dữ liệu, cấu trúc đơn điệu, hàng đợi đơn điệu hoặc tính đơn điệu của quyết
định.

\section{Dãy con tăng dài nhất}

Cho một dãy số nguyên $A$ độ dài $N$. Cần tìm một dãy con dài nhất
$x_1,\ldots,x_k$ sao cho
\[
  1\le x_1<x_2<\cdots<x_k\le N
  \quad\text{và}\quad
  A[x_1]<A[x_2]<\cdots<A[x_k].
\]
Gọi $F[i]$ là độ dài dãy con tăng dài nhất kết thúc tại $A[i]$. Công thức cơ
bản là
\[
  F[i]=\max\{F[j]\mid j<i,\ A[j]<A[i]\}+1.
\]
Thuật toán trực tiếp có $O(N)$ trạng thái và mỗi trạng thái kiểm tra $O(N)$ vị
trí trước đó, nên tổng độ phức tạp là $O(N^2)$.

\subsection*{Tối ưu bằng cấu trúc dữ liệu}

Nút thắt nằm ở phép lấy cực đại với điều kiện $A[j]<A[i]$. Sau khi đã xử lý
các vị trí trước $i$, ta có một tập khóa dạng $(f,x)$, trong đó $f$ là giá trị
quy hoạch động và $x$ là giá trị phần tử tương ứng. Với trạng thái $i$, cần
trả lời truy vấn:
\[
  \max\{f\mid x<A[i]\}.
\]
Sau đó thêm khóa mới $(F[i],A[i])$ vào tập.

Nếu miền giá trị của $A[i]$ có thể nén tọa độ, cây đoạn hoặc cây Fenwick cho
phép thực hiện truy vấn tiền tố và cập nhật điểm trong $O(\log N)$. Khi đó
tổng thời gian giảm từ $O(N^2)$ xuống $O(N\log N)$.

\subsection*{Tối ưu bằng loại bỏ khóa thừa}

Một cách nhìn khác là duy trì chỉ các khóa còn có ích. Nếu tồn tại hai khóa
$(f_1,x_1)$ và $(f_2,x_2)$ sao cho
\[
  f_1\ge f_2,\qquad x_1\le x_2,
\]
thì $(f_2,x_2)$ không bao giờ tốt hơn $(f_1,x_1)$ cho bất kỳ truy vấn
$x<A[i]$ nào, nên có thể bỏ nó. Sau khi loại bỏ toàn bộ khóa bị thống trị,
danh sách còn lại có cả $f$ và $x$ tăng nghiêm ngặt. Truy vấn ``khóa có $x$
nhỏ hơn $A[i]$ và lớn nhất'' có thể trả lời bằng tìm kiếm nhị phân. Khi thêm
khóa mới, ta xóa các khóa bị nó thống trị để khôi phục bất biến.

Cách diễn giải này dẫn đến thuật toán kinh điển duy trì giá trị nhỏ nhất có
thể của phần tử cuối cho mỗi độ dài dãy con tăng.

\section{Dãy con Fibonacci: ZJU 2672}

Một dãy $a_1,a_2,\ldots,a_n$ được gọi là dãy Fibonacci nếu
\[
  a_i=a_{i-1}+a_{i-2}\qquad (i\ge 3).
\]
Với một dãy số nguyên cho trước, bài toán yêu cầu tìm dãy con Fibonacci dài
nhất, với $N\le 3000$.

Gọi $F[i,j]$ là độ dài dãy con Fibonacci dài nhất kết thúc bằng hai phần tử
$A_i,A_j$. Với $i<j$, trạng thái tự nhiên có dạng
\[
  F[i,j]=\max\{F[k,i]+1\mid k<i,\ A_k+A_i=A_j\}.
\]
Số trạng thái là $O(N^2)$, và nếu mỗi trạng thái quét mọi $k$ thì tổng thời
gian là $O(N^3)$.

Điểm quan trọng là với cặp $(i,j)$ đã cố định, ta chỉ cần tìm vị trí $k$ sao
cho $A_k=A_j-A_i$. Nếu cần chọn trong nhiều vị trí có cùng giá trị, vị trí lớn
nhất là lựa chọn có ý nghĩa nhất cho trạng thái kết thúc gần nhất. Có thể tìm
$k$ bằng:
\begin{itemize}
  \item tìm kiếm nhị phân trong danh sách các vị trí của giá trị $A_j-A_i$,
        đạt $O(\log N)$ cho mỗi chuyển;
  \item hoặc duy trì bảng vị trí theo thứ tự quét để truy vấn trong $O(1)$
        trong những biến thể có thể xử lý tuyến tính.
\end{itemize}
Như vậy phần chuyển không còn là phép quét $O(N)$.

\section{Painting the Balls: SGU 183}

Có $N$ quả bóng trên một hàng, ban đầu đều màu trắng. Ta có thể chọn quả bóng
thứ $i$ để tô đen với chi phí $C_i$. Điều kiện là trong mọi đoạn liên tiếp
gồm $M$ quả bóng phải có ít nhất hai quả bóng đen. Hãy tìm tổng chi phí nhỏ
nhất. Ràng buộc điển hình là
\[
  2\le N\le 10000,\qquad 2\le M\le 100,\qquad M\le N.
\]

Trạng thái tự nhiên lưu hai quả bóng đen cuối cùng. Gọi $F[i,j]$ là chi phí
nhỏ nhất khi hai quả bóng đen cuối cùng ở $i$ và $j$ với $i<j$. Khi chọn $j$
làm quả bóng đen mới, quả bóng đen trước $i$ là một vị trí $k$ thỏa điều kiện
không để xuất hiện đoạn dài $M$ chỉ có ít hơn hai quả bóng đen:
\[
  F[i,j]=C_j+\min\{F[k,i]\mid j-M<k<i\}.
\]
Số trạng thái là $O(NM)$, nhưng chuyển trực tiếp mất $O(M)$, cho tổng
$O(NM^2)$.

\subsection*{Nhìn phép chuyển như truy vấn cực tiểu}

Với $i$ và $j$ cố định, quyết định là giá trị nhỏ nhất trên một khoảng chỉ số
$k$. Cây đoạn có thể giảm thời gian quyết định xuống $O(\log M)$.

Tuy nhiên còn có thể tốt hơn. Khi xử lý theo từng giai đoạn $i$ và tính các
$j$ theo thứ tự giảm dần, các khoảng cần lấy cực tiểu trượt dần. Thông tin đã
tính cho $F[i,j+1]$ có thể được tái sử dụng cho $F[i,j]$. Duy trì cực tiểu
của cửa sổ hiện tại bằng hàng đợi đơn điệu giúp mỗi chuyển còn $O(1)$, đưa
tổng thời gian về $O(NM)$.

\section{USACO DEC04 Divide}

Trên một luống hoa độ dài $L$ với $L\le 10^6$, cần lắp các vòi phun. Bán kính
mỗi vòi nằm trong đoạn $[A,B]$, với $1\le A\le B\le 10^3$, và mỗi vị trí phải
được đúng một vòi che phủ. Ngoài ra có $N\le 10^3$ đoạn cho trước; mọi bông
hoa trong cùng một đoạn phải do cùng một vòi che phủ. Mục tiêu là dùng ít vòi
nhất.

Gọi $F[i]$ là số vòi ít nhất khi vòi cuối cùng kết thúc ở vị trí $i$. Khi đặt
vòi cuối che đoạn từ sau $j$ đến $i$, độ dài đường kính phải nằm trong khoảng
$[2A,2B]$, đồng thời điểm cắt $j$ không được nằm bên trong một đoạn bị ràng
buộc. Công thức là
\[
  F[i]=1+\min\{F[j]\mid 2A\le i-j\le 2B,\ j
  \text{ không nằm trong đoạn cấm}\}.
\]
Nếu $j$ nằm trong một đoạn cho trước, ta có thể đặt $F[j]=+\infty$. Khi đó
công thức trở thành truy vấn cực tiểu cửa sổ:
\[
  F[i]=1+\min\{F[j]\mid 2A\le i-j\le 2B\}.
\]
Cây đoạn cho thời gian chuyển $O(\log L)$. Nhưng do cửa sổ chỉ trượt khi $i$
tăng, ta có thể duy trì các ứng viên bằng hàng đợi đơn điệu: loại bỏ phần tử
ra khỏi cửa sổ, và loại bỏ từ cuối hàng đợi các vị trí có giá trị $F$ không
tốt hơn vị trí mới. Khi đó mỗi vị trí vào và ra hàng đợi nhiều nhất một lần,
cho tổng thời gian $O(L)$.

\section{Bất đẳng thức tứ giác và đơn điệu quyết định}

Xét các quy hoạch động hai chiều có trạng thái $F[i,j]$, trong đó quyết định
$k$ chia đoạn và phụ thuộc vào hai phần $F[i,k]$, $F[k,j]$. Ký hiệu quyết định
tối ưu là $K[i,j]$. Nếu có thể chứng minh tính đơn điệu
\[
  K[i,j-1]\le K[i,j]\le K[i+1,j],
\]
thì khi tính $F[i,j]$ ta không cần thử mọi $k$ trong đoạn $[i,j)$. Tổng số lần
liệt kê quyết định trên toàn bộ các trạng thái độ dài đoạn có thể giảm xuống
$O(N^2)$.

\subsection*{Ví dụ: gộp đá}

Bài toán gộp đá có dạng chuyển quen thuộc:
\[
  F[i,j]=\operatorname{sum}(i,j)+
  \min_{i\le k<j}\{F[i,k]+F[k+1,j]\}.
\]
Cách tính trực tiếp thử mọi $k$ cho mỗi cặp $(i,j)$, nên mất $O(N^3)$. Khi
hàm chi phí thỏa điều kiện để quyết định tối ưu đơn điệu, ta chỉ cần xét
\[
  K[i,j-1]\le k < K[i+1,j],
\]
và nhận được lời giải $O(N^2)$.

\section*{Kết luận}

Các ví dụ trên đều bắt đầu từ cùng một câu hỏi: phép chuyển đang thật sự làm
gì? Nếu nó chỉ là lấy cực đại hoặc cực tiểu trên một miền có cấu trúc, ta nên
tìm cấu trúc dữ liệu hoặc bất biến đơn điệu phù hợp. Nếu quyết định tối ưu
dịch chuyển có quy luật, hãy chứng minh tính đơn điệu của quyết định và thu
hẹp phạm vi thử. Đây là tinh thần chính của tối ưu hóa quy hoạch động: không
nhất thiết thay mô hình trạng thái, mà làm cho phép chuyển rẻ hơn.

\end{document}
